\section{The Autumn Budget 2026 choice}
\label{sec:policy}

The results in Section~\ref{sec:results} bear directly on a decision that has not yet been taken. The Chancellor has ruled out universal support of the 2022 kind and signalled income-based help --- in practice, a social tariff --- while the Joseph Rowntree Foundation and the TUC have pressed for a universal discounted block costing around \pounds 5 billion. This section sets out what the two designs do differently, why the disagreement is not really about generosity, and what the model says about the third option nobody is arguing for.

\subsection{Two designs, one disagreement}

A \emph{means-tested social tariff} conditions a discounted unit rate on receipt of a qualifying benefit. Its appeal is arithmetic: every pound of spend goes to a household that has already passed a means test, so on the two conventional targeting metrics --- share of spend reaching the bottom three deciles, and cost per pound of bottom-decile gain --- it will dominate any untargeted alternative, and Table~\ref{tab:scorecard} confirms it does.

The \emph{universal discounted block} does something structurally different. It prices the first half of typical consumption at a discounted rate and everything above it at the full rate, with a per-child allowance recognising that household need scales with size. Because the discount applies to a fixed quantity rather than to all consumption, its cash value is roughly flat across the distribution while high-consumption households --- disproportionately affluent, per \citet{fetzer2024} --- pay the full marginal rate on the bulk of their usage. It is therefore neither a universal price subsidy of the Energy Price Guarantee type, which cost around \pounds 23 billion and was regressive in cash terms precisely because it subsidised every unit, nor a transfer. JRF costs it at roughly \pounds 5 billion and argues it fully offsets the modelled cost increase for deciles one to three. They also argue, on operational rather than distributional grounds, that a social tariff cannot be built in time for winter 2026: there is no register of eligible households that is both current and complete, and constructing one is a multi-year data-sharing project.

The disagreement is not about how much to spend. It is about a coverage gap.

\subsection{The crux: the households outside the means test}

Approximately 40 per cent of households struggling to heat their homes are not on means-tested benefits. This single statistic is the anchor of every submission to the Budget, and it is what our uncompensated-loser metric measures directly. A means test that is perfectly accurate on the population it covers still leaves the entire uncovered population at zero compensation, and that population is not a random sample of the non-poor: it is concentrated among low-income households above the benefit taper, households with savings above the capital limit, pensioners on modest occupational incomes, the self-employed with volatile earnings, and non-claimants --- take-up of means-tested support is materially below entitlement and is lowest among precisely the groups whose energy exposure is highest.

The reform to the Warm Home Discount that took effect for the 2026--27 scheme year --- the qualifying date was 23 August 2026, the property-cost test is scrapped, the \pounds 150 payment is now automatic in England and Wales, and around 2.7 million additional households are brought in --- narrows this gap but does not close it, because eligibility still runs through benefit receipt. The expansion moves the uncompensated share in the right direction; it does not change the structure that generates it.

Figure~\ref{fig:uncompensated} is therefore the policy-relevant figure rather than Table~\ref{tab:scorecard}. Read on between-decile targeting, the social tariff wins. Read on within-decile coverage, it does not, because the losers it misses are systematically the ones no observable identifies. This is \citet{sallee2019}'s result in applied form: when transfers must condition on observables, compensating every loser is not merely expensive but infeasible, and the residual is not noise. \citet{cronin2019} and \citet{douenne2020} establish the same point from the incidence side --- horizontal redistribution within income groups routinely exceeds vertical redistribution between them --- which means a scorecard that reports only decile means is measuring the smaller of the two effects.

\subsection{Against the ``never subsidise prices'' consensus}

The grey-literature consensus after 2022 is unambiguous: compensate incomes, not prices. \citet{ari2022} is the standard reference, and the Bruegel fiscal tracker's finding that only around a third of European energy support was targeted is the standard evidence. The reasoning is familiar --- a price subsidy distorts the marginal incentive to conserve, its cost scales with consumption and therefore with affluence, and it blunts the very price signal that is supposed to induce adjustment.

The peer-reviewed frontier does not agree. \citet{levell2025}, the closest thing to a definitive treatment of the 2022--23 UK episode, estimate a mean welfare loss of about 6 per cent of income absent policy, find that the actual UK package reduced both the mean and the dispersion of losses at an efficiency cost of about 12 per cent of its revenue cost, and --- the result that matters here --- show that the constrained-optimal policy includes a \emph{strictly positive price subsidy}. The subsidy shrinks as the transfer instrument is allowed to condition on more information, in particular on income \emph{and past energy usage} jointly, but it does not go to zero. The reason is exactly the coverage gap of the previous subsection: a price subsidy reaches households that no observable-based transfer can identify, and when the untargeted population is large, the insurance value of that reach outweighs the deadweight loss from the distorted margin.

This reframes the Budget choice. The question is not whether to subsidise prices --- the answer is that a positive subsidy belongs in the optimum --- but how large the price component should be given the transfer instruments actually available. The universal discounted block is a coherent answer to that question and not merely a political compromise: it is a price subsidy on an inframarginal quantity, which preserves the marginal price signal above the block while delivering the reach that a means test cannot. That is close to the two-tier tariff \citet{fetzer2024} propose on entirely separate evidence, and close to the structure \citet{levell2025} recover as optimal when the transfer instrument can condition on past usage --- which, as it happens, the block design does implicitly by defining the block relative to typical consumption.

\subsection{Financing}

The financing side is scored separately, because netting it into the household distribution would conflate two distinct incidence questions.

IPPR's proposal is to claw back windfalls accruing to regulated network companies --- whose allowed returns were set against a very different cost of capital and whose revenues are indexed --- and return the proceeds as a flat household rebate of around \pounds 183. As a distributional object a flat rebate is progressive in percentage terms and reaches every household regardless of benefit status, so it scores well on the coverage metric and poorly on the targeting metrics; it is the mirror image of the social tariff, and Table~\ref{tab:scorecard} shows the two failing on opposite metrics. Its distinctive feature is that its financing is close to a pure rent tax and therefore carries little efficiency cost, which is not true of the alternatives.

The other financing routes on the table are the successor arrangements to the Energy Profits Levy on oil and gas production, and an increase in the Electricity Generator Levy from 45 to 55 per cent. Both are windfall instruments whose yield is mechanically correlated with the shock being compensated, which is an attractive property: the fiscal cost of support and the revenue that pays for it move together, so the package is closer to self-financing in the adverse scenario than in the baseline. Both also carry investment-margin costs that a static microsimulation cannot evaluate and that the source-side literature \citep{goulder2019} suggests are progressive in incidence --- reducing returns to capital falls on capital owners, who are concentrated at the top --- so excluding them from the household results understates the progressivity of the financed package as a whole. I report the household side of each policy and the aggregate cost, and note the direction of the omitted source-side effect rather than attempting to quantify it.
